
\section{Data}\label{sec:data}

The empirical design is a monthly, point-in-time allocation study over listed equity options
and VIX options. The investable equity-option universe is restricted to AAPL, AMZN, GOOGL,
JPM, META, MSFT, NVDA, and TSLA, using OPRA/Databento-derived option panels after basic
mark, expiry, Greek, and liquidity filters. The VIX-option universe is built from raw VIX
chain shards and kept separate because VIX options are volatility derivatives, not options
on an equity spot price. Table~\ref{tab:data_summary} records the sample size, observation
count, train/test split, settlement status, and post-cost layer used in the baseline run.

\begin{table}[H]
\centering
\scriptsize
\IfFileExists{tables/data_summary.tex}
  {\resizebox{0.94\textwidth}{!}{\begin{tabular}{ll}
\toprule
Item & Value \\
\midrule
Raw OPRA-derived equity option rows & 160,315 \\
Raw VIX option rows after filters & 103,650 \\
Monthly snapshot dates & 129 \\
Training dates & 69 \\
Test dates & 60 \\
Primary equity underlyings & AAPL, AMZN, GOOGL, JPM, META, MSFT, NVDA, TSLA \\
VIX option treatment & VX-forward Black-76 Greeks; VIX/VVIX state only \\
VIX settlement source & exact VRO/SOQ (536 rows) \\
VIX headline status & exact VRO/SOQ settlement (excluded expiries disclosed in Section 2) \\
Post-cost layers & executable-cost scenarios (mid, half, full spread) and conservative stress-cost layer \\
Broker execution status & not broker-executed live evidence \\
Rolling option bucket assets & 54 \\
Train/test split & through 2020-12-31 / after 2020-12-31 \\
\bottomrule
\end{tabular}
}}
  {\begin{tabular}{ll}\toprule Object & Implementation \\\midrule
  Sample & generated by the empirical pipeline \\
  Equity options & OPRA/Databento-derived single-name option panels \\
  VIX options & raw VIX chain shards with VX-forward alignment \\\bottomrule\end{tabular}}
\caption{Empirical design summary. The table identifies the option universe, monthly
snapshots, train/test split, VIX settlement status, and post-cost conventions used in the
baseline run.}
\label{tab:data_summary}
\end{table}

The train/test split is frozen at December 31, 2020. Training observations through that date
estimate expected option returns, covariance objects, residual risk, factor exposures, and
implementation screens. Test observations after that date are reserved for out-of-sample
portfolio evaluation. The out-of-sample window runs from January 29, 2021 to April 30,
2026: 60 monthly snapshots spanning 64 calendar months. Four calendar months (May 2021,
March 2022, March 2024, and March 2025) have no ledger row because no representative
contract passed the eligibility screens at those decision dates; they are absent from the
return series rather than imputed. Months in which the book holds only cash collateral are
recorded at the zero financing rate, that is, as exact zero returns under the
zero-financing numeraire of Remark~\ref{rem:zero_financing}; the counterfactual
hurdle/no-trade ledger (\path{artifacts/no_trade_periods.csv}) additionally records, for
each snapshot, the expected-return hurdles at which the optimizer would have stood aside
entirely. Rebalancing is monthly. At each decision date, the research implementation sees
only the option panel, state variables, underlying prices, VX curve, forecast inputs, and
cost inputs observable at or before that date. The realized option return is then measured
at the next monthly realization date or at listed expiry, according to the holding ledger.

Equity-option rolling buckets convert expiring contracts into stable empirical objects
without pretending that options are perpetual assets. For each underlying, right,
moneyness, tenor, and liquidity bucket, the decision date selects a representative listed
contract using contemporaneous marks and filters. The selected option can be a call or a
put, and the optimizer can hold it long or short subject to gross NAV, concentration,
Greek, beta, short-premium, liquidity, stress, and assignment-risk constraints. The return
includes the premium paid or received: a long option that expires worthless loses its
premium weight, while a short option that finishes in the money pays the realized payoff
through NAV accounting.

The VIX-option construction follows a different convention. Raw VIX chain shards are
stacked, OSI symbols are parsed, duplicate \((\text{trade date},\text{symbol})\) rows are
coalesced, and each option is aligned to the relevant VX futures forward. Black--76 Greeks
are computed with respect to the VX forward rather than spot VIX. VIX, VVIX, VIX9D, and
VIX3M enter only as state variables. Exact VRO/SOQ settlement is a claim gate: a VIX option
becomes headline-grade only when the expiry row uses \texttt{vro\_soq\_exact}. Proxy rows,
if any, remain appendix diagnostics and cannot support listed-settlement claims.

The settlement gate has a selection cost that should be stated plainly. Of the 630
representative VIX option rows selected by the monthly screens, 536 (about 85\%) carry a
parsed exact VRO/SOQ settlement and enter the holding ledger; the remaining 94 rows (about
15\%) lack a parsed exact settlement and are \emph{excluded} from settlement-gated
evidence rather than proxied with the VIX close. The exclusions are concentrated in
2015--2018---including 20 required 2017 expiries whose SOQ files could not be parsed
exactly (Table~\ref{tab:vix_download_audit_app})---so the retained VIX rows
under-represent an eventful era for the VIX complex, and the training-window VIX evidence
carries a potential coverage bias in the direction of calmer settlement prints. The
expiry-level exact-versus-proxy comparison is written to the machine-readable audit
\path{artifacts/vix_settlement_audit.csv}.

\begin{table}[H]
\centering
\scriptsize
\IfFileExists{tables/vix_settlement_coverage.tex}
  {\resizebox{0.62\textwidth}{!}{\begin{tabular}{lll}
\toprule
Settlement source & Rows & Headline eligible \\
\midrule
vro\_soq\_exact & 536 & yes \\
\bottomrule
\end{tabular}
}}
  {\begin{tabular}{lll}\toprule Source & Rows & Claim status \\\midrule
  VIX close proxy & generated if VRO absent & diagnostic only \\\bottomrule\end{tabular}}
\caption{VIX settlement coverage. VIX option evidence supports listed-settlement claims
only when every VIX expiry row used in the main tables has \texttt{vro\_soq\_exact}
settlement.}
\label{tab:vix_settlement}
\end{table}

Marks, payoffs, settlement, and costs are treated as first-order design choices. Option
weights are premium exposures per dollar of NAV. The mark source, payoff date, settlement
source, Greek model, state snapshot, and train-end date are recorded before performance is
reported. Post-cost evidence is produced at two deliberately different intensities. The
\emph{executable-cost scenarios} charge entry frictions only: midpoint, half-spread, and
full-spread quoted-spread crossing plus explicit fees, with no-fill rules and capacity
limits. The \emph{conservative stress-cost layer} additionally charges round-trip spread
crossing, slippage, borrow proxies, capacity penalties, simulated margin/stress-capital
funding, and assignment-risk penalties; it is a stress bound rather than an execution
estimate. Section~\ref{sec:results} reconciles the two layers, and
Table~\ref{tab:option_accounting_contract} summarizes the accounting contract each
instrument must satisfy. These conventions are adequate for a research paper on
allocation mechanics because they test whether the option-only Markowitz object can be
constructed, audited, stress-screened, and evaluated out of sample. They are not adequate
for production trading claims because they do not contain broker-routed orders, live quote
depth, fill probabilities calibrated from orders, live margin previews, or broker
reconciliation.

\begin{table}[H]
\centering
\scriptsize
\begin{tabularx}{0.96\textwidth}{L{1.25in}L{1.85in}X}
\toprule
Object & Accounting rule & Main failure if omitted \\
\midrule
Equity option & Premium weight times mark or payoff return & Treats the option as free convexity \\
Short option & Premium received plus future liability & Hides tail and assignment exposure \\
VIX option & Black--76 exposure to matched VX forward & Treats spot VIX as tradable \\
Cash/collateral & NAV numeraire, not optimized asset & Converts the model into a stock-cash allocation \\
Executable-cost scenarios & Mid/half/full quoted-spread entry frictions plus fees & Overstates net returns if spreads are crossed in full \\
Stress-cost layer & Round-trip spreads, fees, slippage, borrow, capacity, margin funding, assignment flags & Promotes gross research returns to false trading evidence \\
\bottomrule
\end{tabularx}
\caption{Option accounting contract. Each object is represented by the cashflow convention
that must be specified before the optimizer can compare it with another option. The two
post-cost rows are distinct layers: the executable-cost scenarios are the realistic
research estimate, the stress-cost layer is a conservative bound
(Section~\ref{sec:results}).}
\label{tab:option_accounting_contract}
\end{table}
